\cotefile{157-3}
\scanpage{177}
\mspage{126}
commutatif dans $\mathrm{Hot}_{W}$ (mais non dans \unsure{$\mathit{Cat}$}, sauf si $X_0=\varnothing$). Supposons que $W$ satisfait $W(7\,\mathrm{bis})$\inserted{(axiome du carré cocartésien, cf p.~75)}. Alors on a
\[
\left\{\begin{aligned}
\alpha_1\in W&\Longleftrightarrow i_2\in W\\
\alpha_2\in W&\Longleftrightarrow i_1\in W .
\end{aligned}\right.
\]
% slanted
\margin{\uline{NB} Pour l'implication $\deleted{\unsure{\alpha}}\Rightarrow$ on n'a besoin que de la forme \unsure{douce A|douce} de $W(7\,\mathrm{bis})$, (i.e. \unsure{(avec} les notations de \uline{XVI}) de $\mathrm{Loc}\,(6\,\mathrm{bis})$}

\deleted{\unsure{Comme}} \uline{Dém} Utilisant le lemme précédent au cas du \unsure{cas|cor.}, on voit que chacune des flèches $\alpha_1,\alpha_2,i_1,i_2$ est dans $W$, ssi la flèche correspondante $\alpha'_1,\alpha'_2,i'_{\unsure{1|2}}$ ou $i'_2$ l'est. Donc les équivalences ci-dessus équivalent à
\[
\begin{aligned}
\alpha'_1\in W&\Longleftrightarrow i'_2\in W\\
\alpha'_2\in W&\Longleftrightarrow i'_1\in W ,
\end{aligned}
\]
lesquelles sont vraies, en vertu de $W(7\,\mathrm{bis})$ \inserted{(formulation \unsure{duale|dure})} appliqué \deleted{\unsure{aux}} à $\mathscr X$ et aux deux \deleted{\unsure{ouverts}}\uline{\unsure{fermés|souscatég.|segments}} $\,{}_{1/}\mathscr X_{\deleted{\unsure{/1}}}$, $\,{}_{2/}\mathscr X_{\deleted{\unsure{/2}}}$ de $\mathscr X$.

\uline{Corollaire} : Considérons \uline{deux} diagrammes de type $\Phi^0$, et un hom. \inserted{$f=(f_0,f_1,f_2)$} de l'un dans l'autre (cf diagramme à traits pleins ci-contre) :
\diagram{\begin{tikzcd}[ampersand replacement=\&,column sep=small,row sep=small]
\& X_1 \arrow[dddd,"f_1" near end,crossing over] \arrow[rrd,dashed,"i_1"] \& \& \\
X_0 \arrow[ru,"\alpha_1"] \arrow[rrd,"\alpha_2"] \arrow[dddd,"f_0"'] \& \& \& \mathscr X \arrow[dddd,dashed,"f\deleted{\unsure{\mathscr X}}"]\\
\& \& X_2 \arrow[ru,dashed,"i_2"] \arrow[dddd,"f_2"'] \& \\
\& \& \& \\
\& X'_1 \arrow[rrd,dashed,"i'_1"] \& \& \\
X'_0 \arrow[ru,"\alpha'_1"] \arrow[rrd,"\alpha'_2"] \& \& \& \mathscr X'\\
\& \& X'_2 \arrow[ru,dashed,"i'_2"] \&
\end{tikzcd}}
d'où un hom
% archive sheet number: 177
